% =====================================================================
%  CAPÍTULO 2 · FUNDAMENTOS Y ESTADO DEL ARTE
%  Función: dar al lector TODO lo que necesita para entender el trabajo y
%  situarlo frente a lo que ya existe. Es donde más pesa el rigor (citas).
%  Reglas clave:
%   - Cada concepto: qué es, cómo funciona, un elemento visual y qué se hace
%     hoy con él. Explica los términos técnicos la primera vez.
%   - Compara (tabla con criterios), no enumeres tecnologías.
%   - Fuentes primarias u oficiales: artículos, documentación, normas.
%     Evita Wikipedia, blogs, vídeos o foros como fuente principal.
%   - No repitas aquí el stack: las herramientas van en el Capítulo 3.
%   - Si hay «Marco teórico» y «Estado del arte» que se repiten, fusiónalos.
% =====================================================================
\chapter{Fundamentos y estado del arte}
\label{cap:estado}

Este capítulo introduce los conceptos necesarios para seguir el resto de la
memoria y analiza las soluciones existentes, con el fin de identificar el hueco
que cubre este trabajo.

\section{Conceptos fundamentales}
\label{sec:conceptos}

\subsection{Arquitectura cliente-servidor}
\label{sec:cliente-servidor}

Presenta el concepto con tus palabras y apóyalo en fuentes. En una arquitectura
cliente-servidor, los clientes envían peticiones a un servidor que centraliza
los datos y la lógica de negocio~\cite{martin2022arquitectura}. La
Figura~\ref{fig:cliente-servidor} muestra el esquema que sigue este trabajo: la
aplicación móvil (izquierda) se comunica mediante una \ac{API} \ac{REST} con el
servidor (centro), que es el único componente que accede a la base de datos
(derecha). Esta separación permite cambiar el cliente sin modificar los datos.

% GUÍA: cita la figura ANTES de que aparezca y guía al lector por ella
% (qué mirar y qué concluir). Un pie solo no basta.
\begin{figure}[htbp]
  \centering
  \begin{tikzpicture}[node distance=2.2cm,
      caja/.style={draw, rounded corners, minimum width=2.8cm, minimum height=1.1cm,
                   align=center, fill=gray!8},
      flecha/.style={-{Stealth[length=2.5mm]}, thick}]
    \node[caja] (cli) {Aplicación\\móvil};
    \node[caja, right=of cli] (srv) {Servidor\\(API REST)};
    \node[caja, right=of srv, shape=cylinder, shape border rotate=90, aspect=0.25,
          minimum width=2.2cm] (db) {Base de\\datos};
    \draw[flecha] (cli.20) -- node[above, font=\small] {petición} (srv.160);
    \draw[flecha] (srv.200) -- node[below, font=\small] {respuesta} (cli.340);
    \draw[flecha, <->] (srv) -- node[above, font=\small] {SQL} (db);
  \end{tikzpicture}
  \caption{Esquema de la arquitectura cliente-servidor del sistema.}
  \label{fig:cliente-servidor}
  \fuente{Elaboración propia a partir de~\cite{martin2022arquitectura}.}
\end{figure}

\subsection{Usabilidad y su medida}
\label{sec:usabilidad}

Cuando un trabajo usa una métrica, defínela con precisión. La \ac{SUS} es un
cuestionario de diez ítems que produce una puntuación entre 0 y
100 y que propuso Brooke~\cite{brooke1996sus}. Para cada participante, la
puntuación se calcula como
\begin{equation}
  \mathrm{SUS} = 2{,}5 \left( \sum_{i\ \text{impar}} (r_i - 1)
               + \sum_{i\ \text{par}} (5 - r_i) \right),
  \label{eq:sus}
\end{equation}
donde $r_i \in \{1,\dots,5\}$ es la respuesta al ítem $i$. La
Ecuación~\eqref{eq:sus} se usa en el Capítulo~\ref{cap:pruebas} para evaluar el
objetivo~\ref{obj:validacion}.

\section{Soluciones existentes}
\label{sec:soluciones}

En un \ac{TFG} de aplicación, este apartado es el estudio de mercado; en un
trabajo de investigación, la revisión de trabajos relacionados. En ambos casos
se analizan las alternativas con los mismos criterios y se resumen en una tabla.
En este ejemplo se han revisado tres aplicaciones (A, B y C), elegidas por ser
las más descargadas de su categoría en noviembre de
2026~\cite{informe_movil2025}, y un \ac{TFG} previo de la \ac{EPS} que abordó un
problema parecido con tecnologías web~\cite{ferrer2025tfg}.

% GUÍA: datos volátiles (descargas, valoraciones de tienda) siempre con fecha
% de consulta. Cada alternativa con su logo o captura si es una app.

La Tabla~\ref{tab:comparativa} compara las alternativas con los cuatro criterios
que se derivan de la motivación. Ninguna de las soluciones analizadas combina la
ocupación en tiempo real con la reserva desde el móvil, que es precisamente el
hueco que pretende cubrir \emph{SalaLibre}.

\begin{table}[htbp]
  \centering
  \caption{Comparativa de las soluciones existentes.}
  \label{tab:comparativa}
  \begin{tabularx}{\textwidth}{@{}lCCCC@{}}
    \toprule
    \textbf{Solución} & \textbf{Ocupación en tiempo real} & \textbf{Reserva móvil}
      & \textbf{Avisos} & \textbf{Código abierto} \\
    \midrule
    Aplicación A~\cite{ejemplo_appA} & $\times$ & \checkmark & \checkmark & $\times$ \\
    Aplicación B~\cite{ejemplo_appB} &\checkmark & $\times$ & $\times$ & $\times$ \\
    Aplicación C~\cite{garcia2024movil} & $\times$ & \checkmark & $\times$ & \checkmark \\
    \midrule
    \emph{SalaLibre} (este trabajo) & \checkmark & \checkmark & \checkmark & \checkmark \\
    \bottomrule
  \end{tabularx}
  \fuente{Elaboración propia a partir de las fuentes citadas (consulta: noviembre de 2026).}
\end{table}

\section{Conclusiones del capítulo}
\label{sec:conclusiones-estado}

Cierra el capítulo con dos o tres frases: qué se ha aprendido y qué decisión
se toma a partir de ello. En este caso, la falta de soluciones que integren
ocupación y reserva justifica el desarrollo de una aplicación propia, cuyos
requisitos se derivan de los criterios de la Tabla~\ref{tab:comparativa}.
